\newcolumntype{C}[1]{>{\centering\arraybackslash}m{#1}}
\chapter{Analyse et spécification des besoins} 
\minitoc
\newpage

\section{Introduction}
Dans le chapitre précédent, nous avons analysé le contexte général de notre projet et défini le cadre méthodologique Agile Scrum. Dans ce deuxième chapitre, nous entrons au cœur de l'ingénierie logicielle. Nous allons d’abord identifier avec précision les différents acteurs de notre système, pour ensuite dégager de manière exhaustive les besoins fonctionnels et non fonctionnels de notre projet. Ces besoins seront traduits sous la forme d'un Product Backlog priorisé structuré en 4 Sprints. Enfin, nous procèderons à une modélisation poussée via le diagramme de cas d'utilisation global et présenterons les choix d'architecture de l'application.

\section{Rôles Scrum}
Dans le cadre de ce projet, l'organisation du travail s'appuie sur les rôles fondamentaux de Scrum afin de clarifier les responsabilités et de faciliter le suivi de l'avancement.

\begin{itemize}
    \item \textbf{Product Owner :} il représente la vision du produit, priorise les besoins et valide les fonctionnalités à développer selon leur valeur métier.
    \item \textbf{Scrum Master :} il veille à la bonne application de la démarche Scrum, facilite la communication et aide à lever les obstacles rencontrés durant les sprints.
    \item \textbf{Équipe de développement :} elle assure l'analyse, la conception, le développement, les tests et l'intégration des fonctionnalités de la plateforme GOSPORT.
\end{itemize}

\section{Spécification des besoins}
La spécification des exigences consiste à traduire les objectifs de haut niveau en un ensemble de fonctionnalités ciblées pour les développeurs et utilisateurs.

\subsection{Identification des acteurs}
Nous distinguons trois acteurs principaux disposant de privilèges d'accès spécifiques :
\begin{itemize}
    \item \textbf{L’Administrateur :} Il interagit via la console web Angular pour assurer la gouvernance globale (statistiques KPIs, modération des utilisateurs, gestion de la bibliothèque d'exercices).
    \item \textbf{Le Coach (ou Nutritionniste) :} Professionnel qui gère son profil, crée des programmes d'entraînement et plans nutritionnels, et suit la progression des athlètes.
    \item \textbf{L’Athlète :} Utilisateur final qui recherche des professionnels, suit ses séances d'entraînement interactives, logge ses repas et suit ses métriques de santé.
\end{itemize}

\subsection{Besoins fonctionnels}
\subsubsection{Besoins de l'Administrateur}
\begin{itemize}
    \item S'authentifier sur le tableau de bord web.
    \item Suivre les indicateurs clés (utilisateurs inscrits, abonnements).
    \item Gérer la bibliothèque globale d'exercices.
    \item Gérer les comptes utilisateurs (bannissement/réactivation).
\end{itemize}

\subsubsection{Besoins du Coach}
\begin{itemize}
    \item Configurer son profil public (spécialités, tarifs).
    \item Valider/refuser les demandes de coaching.
    \item Concevoir des programmes sportifs et plans diététiques.
    \item Suivre la conformité de l'athlète via des historiques et messagerie.
\end{itemize}

\subsubsection{Besoins de l'Athlète}
\begin{itemize}
    \item Créer un compte et compléter son onboarding biométrique.
    \item Rechercher un coach ou nutritionniste et envoyer des demandes.
    \item Lancer des séances d'entraînement via le Workout Player.
    \item Renseigner ses repas pour un calcul immédiat des macronutriments.
    \item Échanger en direct par chat ou visioconférence.
\end{itemize}

\subsection{Besoins non fonctionnels}
\begin{itemize}
    \item \textbf{Sécurité :} Mots de passe hachés avec bcrypt, communication protégée par des tokens JWT.
    \item \textbf{Performance :} Temps de chargement réduits, messagerie instantanée réactive avec WebSockets.
    \item \textbf{Ergonomie :} Design adaptatif moderne (Angular/Tailwind CSS), formulaires réactifs robustes.
\end{itemize}

\section{Product Backlog}
Le Product Backlog est découpé en 4 Sprints selon la méthodologie MoSCoW.

\subsection{Technique de priorisation du Product Backlog}
La priorisation du Product Backlog est réalisée selon la méthode MoSCoW, qui permet de classer les besoins selon leur importance dans la réussite du projet.

\begin{itemize}
    \item \textbf{Must have :} fonctionnalités indispensables au fonctionnement minimal de la plateforme.
    \item \textbf{Should have :} fonctionnalités importantes qui améliorent fortement l'expérience utilisateur.
    \item \textbf{Could have :} fonctionnalités utiles, mais pouvant être planifiées après les besoins prioritaires.
    \item \textbf{Won't have :} fonctionnalités non retenues dans le périmètre initial du projet.
\end{itemize}

\subsection{Product Backlog priorisé}

\begin{table}[H]
\centering
\renewcommand{\arraystretch}{1.4}
\resizebox{\textwidth}{!}{%
\begin{tabular}{|>{\centering\arraybackslash}p{1.2cm}|p{3cm}|>{\centering\arraybackslash}p{1cm}|p{2cm}|p{7cm}|>{\centering\arraybackslash}p{1.5cm}|}
\hline
\rowcolor[HTML]{EFEFEF} 
\textbf{Sprint} & \textbf{Épopée (Epic)} & \textbf{ID} & \textbf{Acteur} & \textbf{User Story} & \textbf{Priorité} \\ \hline
1 & \multirow{3}{3cm}{Socle \& Authentification}
& 1.1 & Tous & ETQ utilisateur, je veux créer un compte sécurisé et utiliser Google OAuth. & Must \\ \cline{3-6}
& & 1.2 & Athlète & ETQ athlète, je veux compléter mon profil d'onboarding biométrique. & Must \\ \cline{3-6}
& & 1.3 & Tous & ETQ utilisateur, je veux me connecter avec un token JWT sécurisé. & Must \\ \hline
2 & \multirow{3}{3cm}{Mise en relation \& Communication}
& 2.1 & Athlète & ETQ athlète, je veux rechercher des professionnels et envoyer des demandes. & Must \\ \cline{3-6}
& & 2.2 & Coach & ETQ coach, je veux accepter ou refuser des demandes de coaching. & Must \\ \cline{3-6}
& & 2.3 & Tous & ETQ utilisateur, je veux chatter en temps réel par messagerie instantanée. & Should \\ \hline
3 & \multirow{3}{3cm}{Entraînement}
& 3.1 & Coach & ETQ coach, je veux construire un programme avec jours et exercices (Workout Builder). & Must \\ \cline{3-6}
& & 3.2 & Athlète & ETQ athlète, je veux exécuter ma séance via le Workout Player interactif. & Must \\ \cline{3-6}
& & 3.3 & Athlète & ETQ athlète, je veux consulter l'historique de mes séances effectuées. & Should \\ \hline
4 & \multirow{4}{3cm}{Nutrition \& IA}
& 4.1 & Coach & ETQ coach, je veux bâtir des plans diététiques (Diet Builder). & Should \\ \cline{3-6}
& & 4.2 & Athlète & ETQ athlète, je veux logger mes repas et suivre mes objectifs de macros. & Should \\ \cline{3-6}
& & 4.3 & Coach/Ath & ETQ utilisateur, je veux consulter le copilote IA pour obtenir des conseils. & Could \\ \cline{3-6}
& & 4.4 & Tous & ETQ utilisateur, je veux réserver des séances et lancer un appel visio. & Could \\ \hline
\end{tabular}%
}
\caption{Product Backlog détaillé en 4 Sprints}
\end{table}

\section{Identification des sprints}
Le découpage en sprints permet de transformer progressivement les besoins en incréments fonctionnels testables. Chaque sprint regroupe un ensemble cohérent de user stories autour d'un objectif métier précis.

\begin{table}[H]
\centering
\renewcommand{\arraystretch}{1.4}
\resizebox{\textwidth}{!}{%
\begin{tabular}{|>{\centering\arraybackslash}p{2cm}|p{5cm}|p{8cm}|}
\hline
\rowcolor[HTML]{EFEFEF}
\textbf{Sprint} & \textbf{Objectif principal} & \textbf{Fonctionnalités couvertes} \\ \hline
Sprint 1 & Socle technique et authentification & Inscription, connexion, JWT, Google OAuth, onboarding utilisateur. \\ \hline
Sprint 2 & Mise en relation et communication & Recherche de professionnels, demandes de coaching/nutrition, messagerie instantanée temps réel (Socket.io). \\ \hline
Sprint 3 & Planification et entraînement & Catalogue d'exercices, Workout Builder (programme), assignation et Workout Player interactif. \\ \hline
Sprint 4 & Nutrition et assistance intelligente & Diet Builder, journal alimentaire, objectifs nutritionnels, visioconférence et copilote IA (Gemini). \\ \hline
\end{tabular}%
}
\caption{Partitionnement des user stories par sprint}
\end{table}

\section{Modélisation et Diagrammes Globaux}
Les figures suivantes représentent le cas d'utilisation global et le diagramme de classes global de l'écosystème web.

\begin{figure}[H]
\centering
\includegraphics[width=0.9\textwidth]{Global_Cas_Utilisation.png}
\caption{Diagramme des cas d'utilisation global}
\end{figure}

\begin{figure}[H]
\centering
\includegraphics[width=0.95\textwidth]{Global_Classes.png}
\caption{Diagramme de classes global}
\end{figure}

\section{Architecture Technique 3-Tiers}
La plateforme adopte une architecture physique 3-Tiers robuste, séparant la présentation, la logique applicative et les données.

\begin{figure}[H]
\centering
\includegraphics[width=0.9\textwidth]{Architecture_3_Tiers.png}
\caption{Architecture physique 3-Tiers}
\end{figure}

\section{Technologies Utilisées}
Dans le cadre de la réalisation de la plateforme GOSPORT, nous avons sélectionné un ensemble de technologies modernes et performantes. Afin de répondre aux exigences de scalabilité, d'interactivité et de robustesse, le choix s'est porté sur un écosystème homogène et éprouvé dans l'industrie.

\subsection{Angular (Interface Web)}
Angular est un framework de développement d'applications web monophages (Single Page Applications - SPA) côté client, développé et maintenu par Google. Reposant entièrement sur TypeScript, Angular impose une architecture hautement structurée en composants, modules, services et directives.

Dans le projet GOSPORT, Angular a été choisi pour concevoir le portail d'administration web et la console de gestion destinée aux professionnels (coachs et nutritionnistes). Sa force réside dans son mécanisme d'injection de dépendances qui rend le code modulaire, son système de routage avancé pour la navigation entre les différents modules, et ses formulaires réactifs (Reactive Forms) puissants qui facilitent la saisie et la validation des données d'onboarding et des programmes d'entraînement.

\begin{figure}[H]
\centering

\includegraphics[width=2.5cm]{angular.png}
\caption{Logo du framework Angular}
\end{figure}

\subsection{Flutter (Application Mobile)}
Flutter est un kit de développement logiciel (SDK) d'interface utilisateur open source développé par Google. Il permet de concevoir des applications compilées nativement pour le mobile (iOS et Android), le web et le bureau à partir d'une seule base de code écrite en Dart.

Pour GOSPORT, Flutter a été sélectionné pour réaliser l'application mobile dédiée aux athlètes. Son moteur de rendu graphique performant permet un affichage à 60 ou 120 images par seconde, assurant des animations fluides indispensables lors de l'exécution des séances via le Workout Player. De plus, son architecture basée sur des widgets réutilisables et son système de rafraîchissement à chaud (Hot Reload) ont grandement accéléré les phases de conception et d'intégration de l'interface utilisateur.

\begin{figure}[H]
\centering

\includegraphics[width=2.5cm]{flutter.png}
\caption{Logo du SDK Flutter}
\end{figure}

\subsection{Node.js et Express (API Backend)}
Node.js est un environnement d'exécution JavaScript côté serveur asynchrone et orienté événements, construit sur le moteur V8 de Google Chrome. Express est un framework Web minimaliste et flexible conçu pour structurer les applications basées sur Node.js.

Ils forment le socle du backend de la plateforme GOSPORT. Ce choix se justifie par les performances exceptionnelles du modèle d'E/S non-bloquant de Node.js, idéal pour gérer des requêtes concurrentes massives (notamment pour le chat en temps réel avec WebSockets). De plus, l'adoption de TypeScript et de l'ORM TypeORM permet de maintenir une structure typée et ordonnée, cohérente avec le frontend Angular, simplifiant ainsi la maintenance du code et l'évolution de la base de données.

\begin{figure}[H]
\centering

\includegraphics[width=3cm]{nodejs.png}
\caption{Logo de Node.js}
\end{figure}

\subsection{PostgreSQL (Base de Données)}
PostgreSQL est un système de gestion de base de données relationnelle et objet (SGBDR) libre et très puissant, réputé pour sa fiabilité, sa robustesse et sa conformité aux standards SQL.

Au sein de GOSPORT, il sert de base de données principale pour stocker toutes les données relationnelles : comptes utilisateurs, profils biométriques, relations coach-athlète, programmes d'entraînement, plans diététiques, logs de nutrition, messages de chat et transactions financières. L'utilisation de PostgreSQL garantit le respect de l'intégrité référentielle, le support des transactions ACID, et offre la flexibilité nécessaire pour exécuter des requêtes d'agrégation complexes indispensables à la génération des graphiques de progression de l'athlète et des indicateurs de performance (KPIs) de l'administration.

\begin{figure}[H]
\centering

\includegraphics[width=2.5cm]{postgresql.png}
\caption{Logo de la base de données PostgreSQL}
\end{figure}

\section{Justification des Choix Technologiques}

Pour chaque brique technologique de GOSPORT, nous avons conduit une analyse comparative afin de retenir la solution la plus adaptée aux contraintes fonctionnelles et techniques du projet. Les tableaux suivants synthétisent ces comparaisons.

\subsection{Base de Données : PostgreSQL vs. MySQL vs. MongoDB}

\begin{table}[H]
\centering
\renewcommand{\arraystretch}{1.4}
\resizebox{\textwidth}{!}{%
\begin{tabular}{|p{5cm}|>{\centering\arraybackslash}p{3.2cm}|>{\centering\arraybackslash}p{3.2cm}|>{\centering\arraybackslash}p{3.8cm}|}
\hline
\rowcolor[HTML]{EFEFEF}
\textbf{Critère} & \textbf{PostgreSQL} & \textbf{MySQL} & \textbf{MongoDB (NoSQL)} \\ \hline
Conformité ACID totale & \checkmark{} Oui & Partielle & \texttimes{} Limitée \\ \hline
Relations complexes (JOIN) & \checkmark{} Excellent & Bon & \texttimes{} Faible \\ \hline
Intégrité référentielle native & \checkmark{} Oui & \checkmark{} Oui & \texttimes{} Manuelle \\ \hline
Support JSON natif & \checkmark{} JSONB natif & Limité & \checkmark{} Natif \\ \hline
Transactions imbriquées (Savepoint) & \checkmark{} Oui & Partielle & \texttimes{} Non \\ \hline
Compatibilité TypeORM & \checkmark{} Excellente & Bonne & Limitée \\ \hline
Performances requêtes analytiques & \checkmark{} Très élevées & Moyennes & Moyennes \\ \hline
\rowcolor[HTML]{E2F0D9}
\textbf{Score d'adéquation GOSPORT} & \textbf{★★★★★} & ★★★☆☆ & ★★☆☆☆ \\ \hline
\end{tabular}%
}
\caption{Comparaison des SGBD -- Justification du choix PostgreSQL}
\end{table}

\textbf{Verdict :} PostgreSQL a été retenu pour ses transactions ACID totales, indispensables lors de la sauvegarde atomique des programmes d'entraînement (cascade \texttt{Program} / \texttt{ProgramDay} / \texttt{ProgramExercise} via QueryRunner). Son support natif du type \texttt{JSONB} est également exploité pour les données nutritionnelles. MongoDB, bien que performant pour les données non structurées, aurait rendu les jointures entre les entités Coach, Athlète, Programme et Repas extrêmement complexes. MySQL, bien que fiable, offre moins de fonctionnalités avancées et une compatibilité TypeORM inférieure.

\subsection{Framework Frontend Web : Angular vs. React vs. Vue.js}

\begin{table}[H]
\centering
\renewcommand{\arraystretch}{1.4}
\resizebox{\textwidth}{!}{%
\begin{tabular}{|p{5cm}|>{\centering\arraybackslash}p{3.2cm}|>{\centering\arraybackslash}p{3.2cm}|>{\centering\arraybackslash}p{3.2cm}|}
\hline
\rowcolor[HTML]{EFEFEF}
\textbf{Critère} & \textbf{Angular} & \textbf{React} & \textbf{Vue.js} \\ \hline
Langage principal & \checkmark{} TypeScript natif & JS / JSX & JS / SFC \\ \hline
Architecture imposée (MVC) & \checkmark{} Oui & \texttimes{} Libre & Semi-imposée \\ \hline
Injection de dépendances native & \checkmark{} Oui & \texttimes{} Externe & \texttimes{} Externe \\ \hline
Formulaires réactifs intégrés & \checkmark{} Oui & Librairies tierces & Librairies tierces \\ \hline
Routage + Guards (AuthGuard, RoleGuard) & \checkmark{} Intégré & Librairies tierces & Librairies tierces \\ \hline
Adapté aux grandes applications & \checkmark{} Excellent & Bon & Moyen \\ \hline
Cohérence avec TypeScript backend & \checkmark{} Totale & Partielle & Partielle \\ \hline
\rowcolor[HTML]{E2F0D9}
\textbf{Score d'adéquation GOSPORT} & \textbf{★★★★★} & ★★★☆☆ & ★★★☆☆ \\ \hline
\end{tabular}%
}
\caption{Comparaison des frameworks frontend -- Justification du choix Angular}
\end{table}

\textbf{Verdict :} Angular a été choisi pour son architecture \textit{opinionated} (stricte et imposée), garantissant une cohérence du code sur l'ensemble du projet. Son système d'injection de dépendances natif, ses Guards de routage (\texttt{AuthGuard}, \texttt{RoleGuard}) et ses Reactive Forms ont permis d'implémenter rapidement les workflows complexes : formulaires d'onboarding, filtrage des coachs, Workout Builder multi-jours. React, bien que très populaire, aurait nécessité l'ajout de nombreuses librairies tierces (Redux, React Router, React Hook Form), augmentant la complexité de maintenance.

\subsection{Environnement d'exécution Backend : Node.js vs. Spring Boot vs. Django}

\begin{table}[H]
\centering
\renewcommand{\arraystretch}{1.4}
\resizebox{\textwidth}{!}{%
\begin{tabular}{|p{5cm}|>{\centering\arraybackslash}p{3.5cm}|>{\centering\arraybackslash}p{3.5cm}|>{\centering\arraybackslash}p{3cm}|}
\hline
\rowcolor[HTML]{EFEFEF}
\textbf{Critère} & \textbf{Node.js + Express} & \textbf{Spring Boot (Java)} & \textbf{Django (Python)} \\ \hline
Modèle d'E/S & \checkmark{} Asynchrone non-bloquant & Synchrone (threads) & Async possible \\ \hline
Support WebSockets (Socket.io) & \checkmark{} Natif & Possible (STOMP) & Possible (Channels) \\ \hline
Partage TypeScript avec Angular & \checkmark{} Totale & \texttimes{} Non & \texttimes{} Non \\ \hline
Consommation mémoire & \checkmark{} Légère ($\sim$50 MB) & Lourde ($\sim$500 MB) & Moyenne \\ \hline
Temps de démarrage & \checkmark{} Très rapide ($<$1s) & Lent (JVM, $\sim$10s) & Rapide \\ \hline
Compatibilité Docker / PaaS & \checkmark{} Excellente & Bonne & Bonne \\ \hline
Requêtes concurrentes (WebSocket) & \checkmark{} Excellent & Bon & Moyen \\ \hline
\rowcolor[HTML]{E2F0D9}
\textbf{Score d'adéquation GOSPORT} & \textbf{★★★★★} & ★★★☆☆ & ★★★★☆ \\ \hline
\end{tabular}%
}
\caption{Comparaison des environnements backend -- Justification du choix Node.js}
\end{table}

\textbf{Verdict :} Node.js a été retenu pour son modèle d'E/S non-bloquant, idéal pour gérer simultanément des centaines de connexions WebSocket (messagerie en temps réel via Socket.io) et de requêtes HTTP concurrentes. L'utilisation de TypeScript -- langue commune avec Angular -- facilite le partage de types entre les couches frontend et backend, réduisant les erreurs d'intégration. Spring Boot, bien que robuste pour les systèmes d'entreprise, implique une empreinte mémoire beaucoup plus lourde, un temps de démarrage de la JVM important et un langage (Java) orthogonal au reste de la stack.

\subsection{Application Mobile : Flutter vs. React Native vs. Xamarin}

\begin{table}[H]
\centering
\renewcommand{\arraystretch}{1.4}
\resizebox{\textwidth}{!}{%
\begin{tabular}{|p{5cm}|>{\centering\arraybackslash}p{3cm}|>{\centering\arraybackslash}p{3.5cm}|>{\centering\arraybackslash}p{3cm}|}
\hline
\rowcolor[HTML]{EFEFEF}
\textbf{Critère} & \textbf{Flutter (Dart)} & \textbf{React Native (JS)} & \textbf{Xamarin (.NET)} \\ \hline
Moteur de rendu & \checkmark{} Propre (Skia/Impeller) & Composants natifs OS & Composants natifs OS \\ \hline
Performance animations & \checkmark{} 60/120 FPS constants & Variable (Bridge JS) & Bonne \\ \hline
Code unique iOS + Android & \checkmark{} Oui & \checkmark{} Oui & \checkmark{} Oui \\ \hline
Hot Reload & \checkmark{} Oui & \checkmark{} Oui & Limité \\ \hline
Richesse des widgets & \checkmark{} Très riche & Bonne & Bonne \\ \hline
Cohérence UI inter-plateformes & \checkmark{} Garantie & Variable & Variable \\ \hline
Tendance communauté & \checkmark{} Forte croissance & Très large & En déclin \\ \hline
\rowcolor[HTML]{E2F0D9}
\textbf{Score d'adéquation GOSPORT} & \textbf{★★★★★} & ★★★★☆ & ★★☆☆☆ \\ \hline
\end{tabular}%
}
\caption{Comparaison des frameworks mobiles -- Justification du choix Flutter}
\end{table}

\textbf{Verdict :} Flutter a été sélectionné pour son moteur de rendu graphique indépendant du système d'exploitation (Skia/Impeller), garantissant des animations constantes à 60 ou 120 images par seconde -- critère essentiel pour le \textbf{Workout Player} de GOSPORT, qui affiche des chronomètres de repos dynamiques et valide des séries en temps réel. Contrairement à React Native, qui délègue le rendu aux composants natifs de l'OS via un pont JavaScript, Flutter garantit une cohérence visuelle parfaite entre iOS et Android à partir d'un unique code source en Dart. Xamarin n'a pas été retenu en raison de son déclin communautaire et de sa complexité de mise en œuvre.

\section{Environnement de Développement}

La qualité et la productivité du développement de GOSPORT reposent sur un ensemble d'outils professionnels complémentaires, formant un environnement de travail cohérent et industriel.

\begin{table}[H]
\centering
\renewcommand{\arraystretch}{1.4}
\resizebox{\textwidth}{!}{%
\begin{tabular}{|p{3.5cm}|p{4cm}|p{8cm}|}
\hline
\rowcolor[HTML]{DAE8FC}
\textbf{Outil} & \textbf{Catégorie} & \textbf{Rôle dans le projet} \\ \hline
Visual Studio Code & Éditeur de code & Développement TypeScript/Angular, extensions ESLint, Prettier, Angular Language Service \\ \hline
Docker Desktop & Conteneurisation & Orchestration des 4 services (API, BDD, UI, pgAdmin) via \texttt{docker-compose.yml} \\ \hline
pgAdmin 4 & Admin base de données & Inspection des tables, exécution de requêtes SQL, supervision des performances \\ \hline
Postman & Test d'API & Collection automatisée avec scripts de test (code HTTP, schéma JSON, temps de réponse) \\ \hline
Git + GitHub & Contrôle de version & Versioning du code source, branches de développement, workflows CI/CD (\texttt{.github/}) \\ \hline
\end{tabular}%
}
\caption{Environnement de développement de GOSPORT}
\end{table}

\subsection{Visual Studio Code}
Visual Studio Code est l'éditeur de code principal utilisé pour le développement de l'ensemble des couches de GOSPORT (backend TypeScript/Express, frontend Angular, scripts utilitaires). Son système d'extensions permet d'enrichir l'environnement avec \textbf{ESLint} et \textbf{Prettier} (qualité du code), l'\textbf{Angular Language Service} (autocomplétion des templates Angular) et \textbf{REST Client} (test rapide des endpoints API directement depuis l'éditeur).

\subsection{Docker et Docker Compose}

\begin{figure}[H]
\centering

\includegraphics[width=2.5cm]{docker.png}
\caption{Logo de Docker}
\end{figure}

Docker est la plateforme de conteneurisation utilisée pour toute l'infrastructure de GOSPORT. Le fichier \texttt{docker-compose.yml} à la racine du projet orchestre les quatre conteneurs de la plateforme en une seule commande (\texttt{docker-compose up -d}). Ces conteneurs communiquent via un réseau bridge dédié (\texttt{coaching-network}) :

\begin{itemize}
    \item \textbf{coaching-app-db :} PostgreSQL 16-alpine, avec healthcheck automatique et volume persistant.
    \item \textbf{coaching-app-api :} Serveur Node.js compilé via un build Docker multi-stage, exposé sur le port 3000.
    \item \textbf{coaching-app-ui :} Dashboard Angular servi par Nginx, exposé sur le port 8081.
    \item \textbf{coaching-app-pgadmin :} Interface pgAdmin4 pour l'administration de la base, sur le port 5050.
\end{itemize}

Ce dispositif garantit un environnement de développement parfaitement reproductible et élimine les problèmes de configuration liés aux dépendances locales (\textit{``works on my machine''}).

\subsection{pgAdmin 4}
pgAdmin est l'interface graphique d'administration de PostgreSQL, intégrée directement dans l'environnement Docker de GOSPORT. Elle permet d'inspecter les tables, d'exécuter des requêtes SQL complexes (agrégations nutritionnelles, vérification des relations), de visualiser les données de test et de diagnostiquer les performances des requêtes lors des phases de débogage. Son déploiement via Docker Compose simplifie l'accès sans installation locale supplémentaire.

\subsection{Postman}
Postman est l'outil de test et de documentation de l'API REST de GOSPORT. Une collection Postman complète regroupe l'ensemble des endpoints de l'API organisés par module : authentification, profils, coaching, programmes, nutrition, messagerie et sessions. Chaque requête est accompagnée de \textbf{scripts de test automatisés} (JavaScript) validant le code HTTP de retour, la structure JSON de la réponse et les contraintes de temps de réponse (ex. : calcul du bilan nutritionnel en moins de 300 ms). Cette collection constitue également la \textbf{documentation vivante} de l'API.

\subsection{Git et GitHub}
Git est le système de contrôle de version décentralisé utilisé pour la gestion du code source de GOSPORT. Le dépôt est hébergé sur GitHub. Le répertoire \texttt{.github/} contient les configurations de workflows d'intégration continue (CI/CD), qui déclenchent automatiquement à chaque \textit{git push} sur la branche principale une vérification de la compilation TypeScript (\texttt{tsc}) et un audit de sécurité des dépendances (\texttt{npm audit}), garantissant l'absence de régressions et de vulnérabilités connues.
